% Canonical Paper II source.
The preceding section fixed the measure and the real energy operator.  We now
complexify the arithmetic frame.  Because $M$ is two-dimensional, the resulting
Cauchy--Riemann equation is the ordinary Riemann-surface equation written in the
assignment-adapted frame; its factorization nevertheless exposes exactly where the
frame bracket and the Laplace--Beltrami drift enter.

\subsection{Compatible complex structure and unitary gauge}

Let $J$ be rotation by $+\pi/2$ in the oriented tangent planes.  Thus
\begin{equation}\label{eq:surface-compatible-complex-structure}
 J\Xu=\Xv,
 \qquad
 J\Xv=-\Xu.
\end{equation}
In real dimension two this almost-complex structure is integrable.  No additional
complex structure is being imposed on the regular AES.

\begin{definition}[Surface arithmetic Cauchy--Riemann operators]
\label{def:surface-arithmetic-cr-operators}
On complex-valued scalar fields define
\begin{equation}\label{eq:surface-arithmetic-cr-operators}
 \dbarAEG:=\frac12(\Xu+i\Xv),
 \qquad
 \dAEG:=\frac12(\Xu-i\Xv).
\end{equation}
These symbols denote first-order scalar coefficient operators.  Intrinsically,
if $(\eta^u,\eta^v)$ is the coframe dual to $(\Xu,\Xv)$, then
\begin{equation}\label{eq:surface-intrinsic-dbar}
 \bar\partial_MF
 =(\dbarAEG F)(\eta^u-i\eta^v).
\end{equation}
\end{definition}

The scalar coefficient in~\eqref{eq:surface-intrinsic-dbar} depends on an
orthonormal frame, while the $(0,1)$-form does not.

\begin{proposition}[$U(1)$-gauge covariance]
\label{prop:surface-u1-gauge}
Let $(Y_u,Y_v)$ be another positively oriented orthonormal frame on an open set
$U\subset M$, related
to the arithmetic frame by a smooth angle $\vartheta$:
\begin{equation}\label{eq:surface-frame-gauge-rotation}
 Y_u=\cos\vartheta\,\Xu+\sin\vartheta\,\Xv,
 \qquad
 Y_v=-\sin\vartheta\,\Xu+\cos\vartheta\,\Xv.
\end{equation}
Then
\begin{equation}\label{eq:surface-cr-gauge-law}
 \frac12(Y_u+iY_v)=e^{-i\vartheta}\dbarAEG.
\end{equation}
Consequently, the equation $\dbarAEG F=0$ is independent of the chosen
oriented orthonormal frame.  On the frame domain $U\subset M$, multiplication on
$L^2(U,d\vol_g)$ by
$e^{-i\vartheta}$ is unitary, so the quadratic operator
$\dbarAEG^*\dbarAEG$ is gauge invariant on compactly supported test fields.
\end{proposition}

\begin{proof}
Expanding~\eqref{eq:surface-frame-gauge-rotation} gives
\[
 Y_u+iY_v
 =e^{-i\vartheta}(\Xu+i\Xv),
\]
which proves~\eqref{eq:surface-cr-gauge-law}.  The dual coframe transforms by
\[
 \theta^u-i\theta^v
 =e^{i\vartheta}(\eta^u-i\eta^v),
\]
so the two gauge factors cancel in~\eqref{eq:surface-intrinsic-dbar}.  Finally,
if $U_\vartheta$ is multiplication by $e^{-i\vartheta}$, then
\[
 (U_\vartheta\dbarAEG)^*(U_\vartheta\dbarAEG)
 =\dbarAEG^*U_\vartheta^*U_\vartheta\dbarAEG
 =\dbarAEG^*\dbarAEG
\]
as a formal-adjoint identity on $C_c^\infty$.
\end{proof}

The arithmetic frame has a more specific gauge relation to the gradient frame
selected by the assignment itself.  This formula identifies exactly which part
of the surface Cauchy--Riemann operator is AEG-specific.

\begin{proposition}[Assignment-specific unitary gauge]
\label{prop:surface-assignment-unitary-gauge}
Let
\begin{equation}\label{eq:surface-gradient-unit-vector}
 q=(\mu^2+\lambda^2a^2)^{1/2},
 \qquad
 E=\frac{\nabla a}{q}.
\end{equation}
Then $(E,JE)$ is a positively oriented orthonormal frame and
\begin{equation}\label{eq:surface-assignment-frame-complex-gauge}
 \boxed{
 \Xu+i\Xv
 =\frac{\mu+i\lambda a}{q}(E+iJE).
 }
\end{equation}
The coefficient
\begin{equation}\label{eq:surface-assignment-u1-factor}
 \chi_a:=\frac{\mu+i\lambda a}{q}
\end{equation}
has unit modulus.  Consequently, if
\[
 \bar Z_E:=\frac12(E+iJE),
\]
then
\begin{equation}\label{eq:surface-assignment-dbar-gauge}
 \dbarAEG=\chi_a\bar Z_E,
 \qquad
 \ker\dbarAEG=\ker\bar Z_E
\end{equation}
pointwise.  Thus the arithmetic Cauchy--Riemann equation has exactly the same
local solution sheaf as the underlying Riemann-surface
$\bar\partial$-equation.
\end{proposition}

\begin{proof}
The eikonal identity~\eqref{eq:paper2-regular-eikonal} gives
$|\nabla a|=q$, so $(E,JE)$ is positively oriented and orthonormal.  The
canonical-frame formula imported from Paper~I is
\[
 \Xu=\frac{\mu}{q}E-\frac{\lambda a}{q}JE,
 \qquad
 \Xv=\frac{\lambda a}{q}E+\frac{\mu}{q}JE.
\]
Adding $i$ times the second identity to the first gives
\begin{align*}
 \Xu+i\Xv
 &=\frac{\mu+i\lambda a}{q}E
   +\frac{-\lambda a+i\mu}{q}JE\\
 &=\frac{\mu+i\lambda a}{q}(E+iJE),
\end{align*}
because $i(\mu+i\lambda a)=i\mu-\lambda a$.  Finally,
\[
 |\chi_a|^2
 =\frac{\mu^2+\lambda^2a^2}{q^2}=1.
\]
The kernel statement follows because $\chi_a$ never vanishes.
\end{proof}

\subsection{Raw factorization and the metric drift}

Let $S_M=\Xu^2+\Xv^2$ as in
\eqref{eq:surface-raw-sum-squares}.  Direct multiplication of the two
first-order expressions gives a factorization which does not yet involve a
measure.

\begin{proposition}[Raw frame factorization]
\label{prop:surface-raw-factorization}
On $C^\infty(M;\C)$, pointwise as differential expressions,
\begin{align}
 4\dAEG\dbarAEG
 &=S_M+i[\Xu,\Xv]
   =S_M+i(c_u\Xu+c_v\Xv),
 \label{eq:surface-raw-factorization-plus}\\
 4\dbarAEG\dAEG
 &=S_M-i[\Xu,\Xv]
   =S_M-i(c_u\Xu+c_v\Xv).
 \label{eq:surface-raw-factorization-minus}
\end{align}
The Laplace--Beltrami expression is instead
\begin{equation}\label{eq:surface-laplacian-drift-comparison}
 \Delta_g=S_M-c_v\Xu+c_u\Xv.
\end{equation}
Thus neither raw product in
\eqref{eq:surface-raw-factorization-plus}--
\eqref{eq:surface-raw-factorization-minus} is, by itself, the metric Laplacian.
\end{proposition}

\begin{proof}
Expand
\[
 (\Xu-i\Xv)(\Xu+i\Xv)
 =\Xu^2+\Xv^2+i(\Xu\Xv-\Xv\Xu)
\]
and reverse the order for the second formula.  Equation
\eqref{eq:surface-laplacian-drift-comparison} is
Proposition~\ref{prop:surface-frame-laplacian}.
\end{proof}

\begin{proposition}[Exact drift factorizations]
\label{prop:surface-exact-drift-factorizations}
The raw products can be written directly in terms of the metric Laplacian:
\begin{align}
 4\dAEG\dbarAEG
 &=\Delta_g+2(c_v+ic_u)\dbarAEG,
 \label{eq:surface-exact-drift-factorization-plus}\\
 4\dbarAEG\dAEG
 &=\Delta_g+2(c_v-ic_u)\dAEG.
 \label{eq:surface-exact-drift-factorization-minus}
\end{align}
These are pointwise identities of differential expressions on
$C^\infty(M;\C)$.
\end{proposition}

\begin{proof}
Using~\eqref{eq:surface-laplacian-drift-comparison},
\begin{align*}
 \Delta_g+2(c_v+ic_u)\dbarAEG
 &=S_M-c_v\Xu+c_u\Xv
   +(c_v+ic_u)(\Xu+i\Xv)\\
 &=S_M+i(c_u\Xu+c_v\Xv)\\
 &=4\dAEG\dbarAEG
\end{align*}
by~\eqref{eq:surface-raw-factorization-plus}.  Complex conjugation of the
coefficients, or the same expansion using
\eqref{eq:surface-raw-factorization-minus}, proves the second identity.
\end{proof}

The distinction between the bracket term and the metric drift is essential.  The
former records the non-coordinate arithmetic frame; the latter is forced by
Riemannian divergence.  They agree in the precise way needed on the kernel of the
Cauchy--Riemann operator, as the next theorem shows.

\subsection{Formal-adjoint factorization}

\begin{theorem}[Adjoint factorization of the surface Laplacian]
\label{thm:surface-adjoint-factorization}
On $C_c^\infty(M^\circ;\C)$, with respect to $d\vol_g$,
\begin{align}
 \dbarAEG^*
 &=-\dAEG+\frac12(c_v+ic_u),
 \label{eq:surface-dbar-formal-adjoint}\\
 \dAEG^*
 &=-\dbarAEG+\frac12(c_v-ic_u).
 \label{eq:surface-d-formal-adjoint}
\end{align}
Consequently,
\begin{equation}\label{eq:surface-adjoint-factorization}
 \boxed{
 4\dbarAEG^*\dbarAEG
 =4\dAEG^*\dAEG
 =-\Delta_g.
 }
\end{equation}
These are formal differential identities on the common test domain; no boundary
condition or self-adjoint domain is implicit in the star notation.
\end{theorem}

\begin{proof}
Using~\eqref{eq:surface-frame-formal-adjoints} and conjugating scalar
coefficients when taking the adjoint,
\begin{align*}
 \dbarAEG^*
 &=\frac12(\Xu^*-i\Xv^*)
   =\frac12(-\Xu+c_v+i\Xv+ic_u)\\
 &=-\dAEG+\frac12(c_v+ic_u),
\end{align*}
which proves~\eqref{eq:surface-dbar-formal-adjoint}; the other formula is its
conjugate counterpart.  Now use
\eqref{eq:surface-raw-factorization-plus}:
\begin{align*}
 4\dbarAEG^*\dbarAEG
 &=-4\dAEG\dbarAEG
   +2(c_v+ic_u)\dbarAEG\\
 &=-S_M-i(c_u\Xu+c_v\Xv)
   +(c_v+ic_u)(\Xu+i\Xv)\\
 &=-S_M+c_v\Xu-c_u\Xv\\
 &=-\Delta_g.
\end{align*}
The calculation for $4\dAEG^*\dAEG$ is analogous.
Appendix~\ref{app:frame-calculations} records both expansions term by term.
\end{proof}

The formal identity has a closed-operator consequence once the domain is stated.

\begin{corollary}[Dirichlet realization by the CR energy]
\label{cor:surface-dirichlet-cr-realization}
Let $\overline{\dbarAEG}_0$ be the graph closure in
$L^2(M,d\vol_g)$ of $\dbarAEG$ initially defined on
$C_c^\infty(M^\circ;\C)$.  Then
\begin{equation}\label{eq:surface-closed-cr-factorization}
 -\Delta_{g,D}
 =4(\overline{\dbarAEG}_0)^*\overline{\dbarAEG}_0.
\end{equation}
The equality is an equality of nonnegative self-adjoint operators represented by
the same closed quadratic form, not merely an equality of pointwise formulas.
\end{corollary}

\begin{proof}
For $f\in C_c^\infty(M^\circ;\C)$,
Theorem~\ref{thm:surface-adjoint-factorization} gives
\begin{equation}\label{eq:surface-cr-energy-equality}
 4\|\dbarAEG f\|_{L^2(d\vol_g)}^2
 =\langle-\Delta_gf,f\rangle
 =\mathcal E_0(f,f).
\end{equation}
Moreover,
\[
 \|f\|_{L^2(d\vol_g)}^2+4\|\dbarAEG f\|_{L^2(d\vol_g)}^2
 =\|f\|_{\mathcal E}^2.
\]
Thus $\dbarAEG$ is closable with graph domain $H_0^1(M,g)$, and the closure
of the form on the left of
\eqref{eq:surface-cr-energy-equality} is exactly the Dirichlet form
$\mathcal E$ from Definition~\ref{def:surface-energy-domain}.  The uniqueness
of the self-adjoint operator represented by a closed nonnegative form proves
\eqref{eq:surface-closed-cr-factorization}.
\end{proof}

\subsection{Arithmetic holomorphicity and harmonicity}

\begin{definition}[Arithmetic holomorphic field on a regular AES]
\label{def:surface-arithmetic-holomorphic}
Let $\Omega\subset M$ be open.  A field $F\in C^1(\Omega;\C)$ is
\emph{arithmetic holomorphic} if
\begin{equation}\label{eq:surface-arithmetic-holomorphic-equation}
 \dbarAEG F=0
 \quad\text{on }\Omega.
\end{equation}
Equivalently, if $F=f+ih$, then
\begin{equation}\label{eq:surface-arithmetic-cr-real-system}
 \Xu f=\Xv h,
 \qquad
 \Xv f=-\Xu h.
\end{equation}
\end{definition}

By Proposition~\ref{prop:surface-u1-gauge}, this definition is independent of
the arithmetic frame used to write the coefficient equation.  By
\eqref{eq:surface-intrinsic-dbar}, it is exactly ordinary holomorphicity for the
Riemann surface $(M,J)$.

\begin{proposition}[Elementary holomorphic calculus]
\label{prop:surface-elementary-holomorphic-calculus}
Let $F=f+ih$ be arithmetic holomorphic on an open set $\Omega\subset M$.
Then the following statements hold.
\begin{enumerate}
 \item If $\Phi$ is classically holomorphic on a complex neighborhood of
 $F(\Omega)$, then $\Phi\circ F$ is arithmetic holomorphic.
 \item The two surface derivative vectors are orthogonal and have equal length:
 \begin{align}
  (\Xu f)(\Xv f)+(\Xu h)(\Xv h)&=0,
  \label{eq:surface-holomorphic-orthogonality}\\
  (\Xu f)^2+(\Xu h)^2
  &=(\Xv f)^2+(\Xv h)^2.
  \label{eq:surface-holomorphic-equal-length}
 \end{align}
 The oriented Jacobian satisfies
 \begin{equation}\label{eq:surface-holomorphic-jacobian}
  J_g(F)
  :=(\Xu f)(\Xv h)-(\Xv f)(\Xu h)
  =(\Xu f)^2+(\Xv f)^2\ge0.
 \end{equation}
 \item If $\Omega$ is connected and $F=H\circ a$ for a
 $C^1$ function $H$ on an interval containing $a(\Omega)$, then $F$
 is constant.
\end{enumerate}
\end{proposition}

\begin{proof}
The real vector fields $(\Xu,\Xv)$ obey the ordinary chain rule.  Since
$\Phi$ is holomorphic, its complex derivative gives
\[
 \dbarAEG(\Phi\circ F)
 =\Phi'(F)\dbarAEG F=0,
\]
which proves the first statement.  The real Cauchy--Riemann system
\eqref{eq:surface-arithmetic-cr-real-system} says
\[
 \Xu f=\Xv h,
 \qquad
 \Xu h=-\Xv f.
\]
Substitution proves
\eqref{eq:surface-holomorphic-orthogonality}--
\eqref{eq:surface-holomorphic-jacobian}.

For the final statement, the assignment derivative rules
\eqref{eq:paper2-assignment-frame-rules} give
\begin{equation}\label{eq:surface-a-only-holomorphic-rigidity}
 2\dbarAEG(H\circ a)
 =(\mu+i\lambda a)H'(a).
\end{equation}
The factor $\mu+i\lambda a$ never vanishes for real $a$, because
$\mu\ne0$.  Hence $H'(a)=0$ throughout $\Omega$, so
$d(H\circ a)=0$.  Connectedness makes $F$ constant.
\end{proof}

\begin{theorem}[Holomorphic fields are Laplace--Beltrami harmonic]
\label{thm:surface-holomorphic-implies-harmonic}
If $F\in C^2(\Omega;\C)$ is arithmetic holomorphic, then
\begin{equation}\label{eq:surface-holomorphic-harmonic-equation}
 \Delta_gF=0
 \quad\text{pointwise on }\Omega.
\end{equation}
Hence the real and imaginary parts of $F$ are real
Laplace--Beltrami harmonic functions.
\end{theorem}

\begin{proof}
The formal-adjoint formula
\eqref{eq:surface-dbar-formal-adjoint} is also an identity of local
differential expressions.  Therefore
\[
 -\Delta_gF
 =4\dbarAEG^*(\dbarAEG F)=0
\]
pointwise.  Taking real and imaginary parts proves the final statement.

Equivalently, the raw factorization gives
\[
 0=4\dAEG\dbarAEG F
   =S_MF+i(c_u\Xu+c_v\Xv)F.
\]
On the Cauchy--Riemann kernel, $\Xu F=-i\Xv F$, and hence
\[
 i(c_u\Xu+c_v\Xv)F
 =(-c_v\Xu+c_u\Xv)F.
\]
The right-hand side converts $S_MF$ exactly into $\Delta_gF$ through
\eqref{eq:surface-laplacian-drift-comparison}.
\end{proof}

\begin{remark}[Local equation versus Hilbert-space kernel]
\label{rem:surface-local-versus-l2-kernel}
Theorem~\ref{thm:surface-holomorphic-implies-harmonic} is local and pointwise;
it does not require $F\in L^2$.  By contrast, a statement such as
$F\in\ker(-\Delta_{g,D})$ requires
$F\in\Dom(-\Delta_{g,D})$ and the
chosen Dirichlet boundary condition.  Many holomorphic fields used later for
ideal-boundary data are not square-integrable with respect to hyperbolic volume,
so the two assertions must not be interchanged.
\end{remark}
